\section{Conclusion} \label{sec:conclusion}

In this paper, we use a cluster randomized controlled trial to evaluate the impacts of a financial literacy pilot program in Brazilian elementary schools during the academic year of 2015. The pilot was implemented in 101 municipal schools and involved approximately 9,000 students. Our main findings showed that the program improved students' financial proficiency, consumption (self-control), and saving (risk-aversion) attitudes, and some behavioral outcomes, such as the use of piggy banks. We find strong indications of heterogeneous effects though. While proficiency gains were stronger among middle school students, changes in attitudes and behavioral outcomes stood out among elementary education students. 

We tested the overlooked hypothesis in this literature that poses that change in formal knowledge precedes changes in actual behavior. We use a causal mediation effect analysis to test this hypothesis. Overall, our results seem to reject the hypothesis that claims that changes in knowledge are a necessary condition to change individuals' habits. 

Finally, we used student test score data in math and reading in different school grades to investigate the program's impacts on contemporary learning outcomes and on human capital accumulation. We do not find evidence that the program impacted grade progression, retention, dropout rates, or learning outcomes. Our reading of the results is that the program did not affect students' intertemporal decision-making, confirming the preliminary results used by the implementing partner who decided to scale the program down in its current format. 

